\documentclass[12pt]{article}
\usepackage[utf8]{inputenc}
\usepackage[T1]{fontenc}
\usepackage[margin=1in]{geometry}
\usepackage{setspace}
\usepackage{xcolor}
\usepackage[colorlinks=true, urlcolor=blue!60!black, linkcolor=blue!60!black]{hyperref}
\setlength{\parindent}{2em}
\setlength{\parskip}{1.25em}
\renewcommand{\baselinestretch}{1.0}
\begin{document}
\begin{center}
{\LARGE \textit{New York v. United States}, 505 U.S. 144 (1992)\par}
\vspace{0.8em}
{\large Case Brief --- Four-Page Edition\par}
\end{center}
\vspace{1em}

\section{Memory Jogger}

Congress may encourage States to regulate through conditional spending or conditional preemption, but the Tenth Amendment's anti-commandeering principle forbids it to compel state legislatures to enact or administer a federal regulatory program, dooming the radioactive-waste Act's ``take title'' provision (\href{https://www.courtlistener.com/opinion/112768/new-york-v-united-states/}{505 U.S. 144}).

\section{Detailed Case Facts}

By 1979 only three commercial sites for low-level radioactive waste remained open (in Nevada, Washington, and South Carolina), and those States began cutting volume and threatening closure. On proposals from the National Governors' Association, Congress made each State responsible for its own waste in 1980 and allowed regional compacts to exclude outside waste. With 31 States still lacking outlets, Congress enacted a state-brokered compromise in the Low-Level Radioactive Waste Policy Amendments Act of 1985: the sited States would accept outside waste through 1992 for graduated surcharges, and unsited States faced three incentives---\textbf{monetary} rebates from escrowed surcharges for meeting siting milestones; \textbf{access} penalties allowing sited States to raise surcharges and eventually deny access; and the \textbf{take-title} provision, requiring any State unable to dispose of its waste by 1996 to take title to and possession of it at the generator's request and assume liability for damages.

New York, a major waste generator, declined to join a compact and sought to site its own facility, but residents of Allegany and Cortland Counties, where five candidate sites lay, fiercely opposed it.

\section{Procedural History}

In 1990 New York and the two counties sued the United States, challenging the Act under the Tenth Amendment, the Guarantee Clause, and other provisions; Washington, Nevada, and South Carolina intervened. The District Court dismissed, the Second Circuit affirmed, and the Supreme Court granted certiorari, deciding the case on June 19, 1992.

\section{Judicial Votes}

\textbf{9--0} to uphold the monetary and access incentives; \textbf{6--3} to invalidate and sever the take-title provision.
\begin{itemize}
\item \textbf{Majority:} Justice O'Connor, joined by Chief Justice Rehnquist and Justices Scalia, Kennedy, Souter, and Thomas, and by Justices White, Blackmun, and Stevens as to Parts III-A and III-B.
\item \textbf{Concurring in part and dissenting in part:} Justice White, joined by Justices Blackmun and Stevens.
\item \textbf{Concurring in part and dissenting in part:} Justice Stevens.
\end{itemize}

\section{Holding}

The monetary and access incentives are valid exercises of the Spending and Commerce powers, but the take-title provision violates the Tenth Amendment and is severable. Congress may not commandeer state legislative processes by directly compelling States to enact and enforce a federal program. It may attach conditions to spending under \textit{South Dakota v. Dole} or offer States a choice between regulating under federal standards and federal preemption, but it may not force a choice between two unconstitutionally coercive alternatives. State officials' consent cannot ratify a departure from the federal structure. The discussion suggesting some Guarantee Clause claims may be justiciable is dicta.

\section{Analysis of Opinions}

\textbf{Majority (O'Connor).} The Tenth Amendment and Article I are ``mirror images of each other.'' Distinguishing \textit{Garcia} as involving generally applicable laws, O'Connor relied on \textit{Hodel}, \textit{FERC}, and the Convention's rejection of the New Jersey Plan to conclude that the Framers chose direct regulation of individuals. Commandeering also blurs political accountability: state officials bear the blame while federal officials escape it. The monetary incentives were valid commerce, tax, and spending measures; the access incentives burdened private generators, not States. But the take-title provision ``crossed the line distinguishing encouragement from coercion'': taking title was a compelled subsidy, and legislating on command was a direct order. Neither a strong federal interest nor New York's earlier support could save it, because federalism protects individual liberty, not officials' convenience.

\textbf{White (with Blackmun and Stevens).} The Act was a state-initiated interstate bargain, and New York, having enjoyed seven extra years of waste exports, should be estopped. \textit{Hodel}'s language was dicta, \textit{Garcia} left federalism to the political process, and the decision perversely encourages Congress to preempt States entirely.

\textbf{Stevens.} The Articles of Confederation let Congress command States, and the Constitution enlarged federal power; Congress routinely directs state functions, so nothing forbids commanding States to implement federal waste standards.

\section{Examples: Future Applications}

\begin{enumerate}
\item \textbf{Same-Side (Commandeering): Carbon-Cap Mandate.} A statute requiring every legislature to cap emissions or else pay residents' pollution damages offers a ``choice'' between two coercive commands, like take title.
\item \textbf{Same-Side (Commandeering): Ban on Repealing Gambling Laws.} Barring legislatures from authorizing sports betting directly orders them, as \textit{Murphy v. NCAA} later held.
\item \textbf{Opposite-Side (Valid): Conditional Preemption.} A Clean Air Act provision letting EPA regulate private factories directly if a State declines to adopt a plan leaves the State free and burdens private emitters.
\item \textbf{Opposite-Side (Valid): Generally Applicable Regulation.} A privacy law barring both data brokers and state motor-vehicle departments from selling drivers' data regulates States as database owners, as in \textit{Reno v. Condon}.
\item \textbf{Fence-Sitter: Duties on State Child-Welfare Agencies.} A statute imposing ``active efforts'' and record-keeping duties on state agencies and courts, as litigated in \textit{Haaland v. Brackeen}, straddles evenhanded regulation, permissible directives to state courts, and executive commandeering.
\end{enumerate}

\section{Critique}

\textbf{Originalist and textualist views.} Defenders praise the revival of structural federalism grounded in the shift from the Articles to direct national legislation. Critics note O'Connor's concession that the limit ``is not derived from the text of the Tenth Amendment itself,'' leading scholars such as Martin Redish to question a categorical bar built on structural inference.

\textbf{Progressive and functionalist views.} Scholars doubt the accountability rationale, since voters often know when States act under federal compulsion and preserved cooperative schemes blur accountability too. Echoing White, they warn that striking a state-negotiated compromise pushes Congress toward full preemption or aggressive spending conditions.

\section{Key Quotations}

\begin{quote}
``In providing for a stronger central government, therefore, the Framers explicitly chose a Constitution that confers upon Congress the power to regulate individuals, not States.''
\end{quote}
\noindent --- Justice O'Connor, Opinion of the Court.

\begin{quote}
``A choice between two unconstitutionally coercive regulatory techniques is no choice at all. Either way, `the Act commandeers the legislative processes of the States by directly compelling them to enact and enforce a federal regulatory program,' \dots{} an outcome that has never been understood to lie within the authority conferred upon Congress by the Constitution.''
\end{quote}
\noindent --- Justice O'Connor, Opinion of the Court.

\begin{quote}
``States are not mere political subdivisions of the United States. State governments are neither regional offices nor administrative agencies of the Federal Government. The positions occupied by state officials appear nowhere on the Federal Government's most detailed organizational chart.''
\end{quote}
\noindent --- Justice O'Connor, Opinion of the Court.

\end{document}
